\cotefile{156-7}
\scanpage{108}
\mspage{107}
pour
\[
\mathrm{omb}(X)^{\unsure{o|0|\circ}}\ne\varnothing\qquad\forall X\in M,
\]
car si on avait $\mathrm{omb}(X)^o=\varnothing$, on aurait $X\mathrel{\unsure{|\circ|\,|\,|o|}}X$, ce qui est exclu par Mag $5''$ (\unsure{antiréflexivité} de $|\circ|$). De plus, elle implique aussi At L2 (via $(***)$) et At L3 (via \deleted{Mag} Cor. de Mag $4''$).

Et l'axiome At C (\unsure{spec|sup}) sera automatiquement valable (alors qu'il ne l'est peut-être pas, et qu'il est \unsure{boiteux|boîteux}, dans le cas où on ne suppose pas $(****)$).

Inversement, si on part d'un $\unsure{\gtrless|\lessgtr}$-magasin, satisfaisant At L1, L2, L3 (pour pouvoir y définir raisonnablement $|\circ|$), la structure $\gtrless$ n'est autre que celle $\gtrless'$ définie par $|\circ|$, ssi $M$ satisfait At C (spec).

Donc le domaine commun des deux notions est, \unsure{résumé} par l'équivalence de catégories :
\diagram{\begin{tikzcd}[ampersand replacement=\&,row sep=large]
\text{$\gtrless$-magasins satisfaisant At L1, L2, L3 et At C (spec)} \arrow[d, Leftrightarrow] \\
\text{$|\circ|$-magasins satisfaisant $(****)$}
\end{tikzcd}}

Donc les $\gtrless$-magasins\inserted{(compatibiliseurs)} et les $|\circ|$-magasins\inserted{(\unsure{disparateurs})} apparaissent comme deux \deleted{types de} généralisations d'une même espèce de structure. Celle-ci paraît trop restrictive à mon sens, par la seule présence de At L1.
\scanpage{109}
\mspage{108}
On dispose donc de deux généralisations possibles, en \unsure{direction} d'une topologie sans axiome de divisibilité. Mais c'est la direction $\unsure{|\circ|\,|\,|o|}$ qui me semble la plus riche. En effet, elle a toujours une structure de $\unsure{\gtrless|\lessgtr}$-\unsure{magasin|magma} correspondante, mais \unsure{avec}, \uline{en plus}, une notion $|\circ|$ de disjonction intrinsèque, qui exprime une intuition géométrique \unsure{palpable}, laquelle échappe (en l'absence de At L1) au seul $\gtrless$-\unsure{point} de vue. L'axiome axiome At C (\unsure{spec|sup}), \unsure{avec} le \unsure{recul}, apparaît artificiel -- il exprime l'embarras de celui qui ne dispose d'une notion de ``disjonction'' en termes des \unsure{illusoires} ``strates ouvertes'' $X^{\unsure{o|0|\circ}}$.

Donc finalement, je change d'optique -- la ribambelle d'axiomes de tout poil At C (spec), At C, At D\ldots\ \unsure{tombe} \unsure{en} un \unsure{bruit} \unsure{floche|flou}, Je reprends
\scanpage{110}
\mspage{109}
tout, encore une fois !

Un \uline{\unsure{magasin}} (de \unsure{multistrates}) est la donnée d'un ensemble $M$, muni de trois relations
\[
\leq,\quad\ll,\quad\unsure{|\circ|\,|\,|o|}
\]
avec la terminologie introduite par ailleurs. La relation $X\mathrel{|\circ|}Y$ se lit ``$X$ et $Y$ intérieurement \unsure{disjoints|disjoint}'', on pourra aussi plus loin l'écrire $X^{\unsure{o|0|\circ}}\mathrel{\|}Y^o$ (quand on aura donné un sens à $X^o$, $Y^o$\ldots).

Les axiomes sont les suivants

% The block from Mag 1 through Mag 4, including its slanted notes, is cancelled by long diagonal strokes.
\deleted{\uline{Mag 1}} \deleted{$X\leq Y$ et $X\ll Y$ sont des relations d'ordre, et}
\[
\deleted{X\leq Y\Longrightarrow X\ll Y}
\]

\deleted{\uline{Mag 2}}
\[
\begin{aligned}
&\deleted{\forall X\in M,}\quad \deleted{\unsure{\forall Z}}\quad \deleted{Z\in M}\quad\deleted{\text{avec }Z\ll X,\ \exists X'\text{ avec}}\\
&\deleted{Z\mathrel{\overset{o}{\ll}}X'\leq X}
\end{aligned}
\]
% slanted
\margin{\deleted{i.e. $\exists$ \unsure{élém\textsuperscript{t}} minimal dans $\widetilde{Y}^{Z}$}}
\deleted{où on écrit}
\[
\begin{aligned}
&\deleted{Z\mathrel{\overset{o}{\ll}}Y\ \overset{\mathrm{def}}{\Longleftrightarrow}\ Z\text{ est élément minimal de}}\\
&\deleted{\widetilde{\unsure{Y|X}}^{\unsure{Z}}=\{Y'\in\widetilde{Y}\mid Z\ll Y'\}}
\end{aligned}
\]

\deleted{\uline{Mag 3}} \deleted{Soient}
\diagram{\begin{tikzcd}[ampersand replacement=\&,row sep=small,column sep=small]
\deleted{Y} \arrow[r,phantom,"\deleted{\ll}"] \& \deleted{X} \arrow[d,phantom,"\deleted{\leq}"{rotate=90}] \\
\& \deleted{X'}
\end{tikzcd}}
\deleted{et $Z$ \deleted{$\unsure{\in}$} \unsure{tels|tel} que $Z\ll Y,X'$, alors $\exists Y'$ avec}
\diagram{\begin{tikzcd}[ampersand replacement=\&,row sep=small,column sep=small]
\& \deleted{Y} \arrow[r,phantom,"\deleted{\ll}"] \arrow[d,phantom,"\deleted{\ll}"{rotate=90}] \& \deleted{X} \arrow[d,phantom,"\deleted{\leq}"{rotate=90}] \\
\deleted{Z} \arrow[r,phantom,"\deleted{\ll}"] \& \deleted{Y'} \arrow[r,phantom,"\deleted{\ll}"] \& \deleted{X'}
\end{tikzcd}}
% slanted
\margin{\deleted{Variante (\unsure{équivalente}) Si $Z\mathrel{\overset{o}{\ll}}Y$, $Z\ll X'$ alors $Y\ll X'$ i.e. $Z$ est un plus \uline{petit} \uline{élément} dans $\unsure{\widetilde{Y}^{Z}}$.}}
\deleted{(i.e.}
\[
\deleted{\mathrm{Omb}(Y)\cap\mathrm{Omb}(X')=\mathrm{Omb}\,Y_{X'}}
\]
\deleted{où $Y_{X'}\overset{\mathrm{def}}{=}\widetilde{Y}\cap\mathrm{Omb}\,X'$)}

\deleted{De plus, si $Y\leq X$, \unsure{OPS|ops} $Y\leq X'$.}

\deleted{\uline{Mag 4}}
\scanpage{111}
\mspage{110}
\uline{Mag 1} $X\leq Y$ et $X\ll Y$ sont des relations d'ordre, et
\[
X\leq Y\Longrightarrow X\ll Y
\]

\uline{Mag 2} \unsure{Soient|Soit} $Z\ll X$, alors l'ensemble
\[
\widetilde{X}^{Z}=\{X'\deleted{\unsure{\in X}}\in\widetilde{X}\mid Z\ll X'\}
\]
a un \uline{plus petit élément}.

Si c'est $X$, on écrit $Z\mathrel{\overset{\unsure{o|0|\circ}}{\ll}}X$

\uline{Mag 3} $X\mathrel{\unsure{|\circ|\,|\,|o|}}Y$ est symétrique et \unsure{antiréflexive}, i.e. elle implique $Y\mathrel{|\circ|}X$ et $X\ne Y$. De plus On a
\[
X\mathrel{|\circ|}Y,\quad X'\mathrel{\overset{o}{\ll}}X,\quad Y'\mathrel{\overset{o}{\ll}}Y\Longrightarrow X'\mathrel{|\circ|}Y'
\]

\uline{Mag 4} Si $X\in M$, $Y,Z\leq X$, $\deleted{\unsure{Y\ll}}$ alors
\[
Y\deleted{\unsure{\leq|\ne}}\ne Z\Longrightarrow Y\mathrel{|\circ|}Z
\]

En termes de \unsure{ces} données, on définit
\[
\begin{aligned}
X\mathrel{\unsure{\gtrless|\lessgtr}}Y\ &\overset{\mathrm{def}}{\Longleftrightarrow}\ \forall X'\in\widetilde{X}\setminus\widetilde{X}\cap\widetilde{Y},\ Y'\in\widetilde{Y}\setminus\widetilde{X}\cap\widetilde{Y},\ \text{on a }X'\mathrel{|\circ|}Y'\\
X\mathrel{\|}Y\ &\overset{\mathrm{def}}{\Longleftrightarrow}\ \forall X'\in\widetilde{X},\ Y'\in\widetilde{Y},\ \text{on a }X'\mathrel{|\circ|}Y'
\end{aligned}
\]
On \deleted{\unsure{vérifie}}\inserted{a} alors
\[
\begin{aligned}
X\mathrel{\|}Y\ &\Longleftrightarrow\ X\mathrel{\gtrless}Y\text{ et }\widetilde{X}\cap\widetilde{Y}=\varnothing\quad\text{(i.e. }\{X,Y\}\\
&\hspace{10em}\text{non minoré dans }M\text{)}
\end{aligned}
\]

\deleted{\uline{\unsure{Filtre}}} \uuline{Figure}\inserted{$F$} (d'un \unsure{magasin}) : partie $\leq$-fermée, telle que $\forall X,Y\in F$, on ait $X\mathrel{\gtrless}Y$.

\deleted{\uline{Atelier} : \unsure{Magasin}\unsure{?}} \uline{Ex} : les $\widetilde{X}$
\scanpage{112}
\mspage{111}
\uline{Atelier}\inserted{(\unsure{spéciaux|speciaux})} : est donné par $(\unsure{\mathcal{F}|\mathscr{F}},\leq,\ll;\unsure{|\circ|\,|\,|o|})$ où $\mathcal{F}$ est un ensemble muni de deux relations d'ordre $\leq$, $\ll$, et d'une \deleted{\unsure{donnée|relation} \unsure{supplémentaire}} relation. $\deleted{\unsure{\gtrless|\lessgtr}}$ $\deleted{\unsure{|\circ|}}$ Axiomes

\uline{At 1} $\mathcal{F}$\inserted{$\ne\varnothing$ et $\mathcal{F}$} stable par $\mathrm{Sup}^{\leq}$ \deleted{de \unsure{f.}} \uline{majorés} quelconques.

\uline{At 2} Toute figure\inserted{$F$} est le $\mathrm{Sup}^{\leq}$ de \inserted{l'ens $\widetilde{F}$ de} ses ``strates''
\[
\widetilde{F}=\{X\in\mathcal{F}\mid X\leq F,\ X\text{ ``strictement irréductible''}\}
\]
% slanted
\margin{\unsure{équivalent à} : $F\leq G\Longleftrightarrow\widetilde{F}\subset\widetilde{G}$ (critère \unsure{des deux} figures par strates) $\widetilde{F}$, \deleted{\unsure{s'app|s'ap}} muni de $\leq$, s'appelle ``la \uuline{\unsure{maquette}} de la figure'' (qu'elle \unsure{restitue|reconstitue} par ``déploiement'').}

\deleted{\unsure{At 3}} Si $M$ désigne l'ens des figures élémentaires -- multistrates (objets str. irréd. de $\mathcal{F}$), $\deleted{\unsure{M}}$ muni de la structure $\leq$ seulement, \unsure{donc|alors} la donnée de $(\mathcal{F},\leq)$ satisfaisant \unsure{:|à} At 1, At 2 \deleted{\unsure{pure}} est équivalente à celle d'un ensemble ordonné $(M,\leq)$, avec un ens.\inserted{$\widetilde{\mathcal{F}}$} de parties fermées\inserted{de $M$} (les $\widetilde{F}$), contenant avec chaque $\unsure{\Phi|\Omega}$ ses parties fermées, \deleted{\unsure{et}} recouvrant $M$ \deleted{\unsure{les}\ $\unsure{\widetilde{F}}$} et contenant au moins la partie vide de $M$ (``figure vide''). Ceci posé
% slanted
\margin{Les $\mathrm{Sup}$\inserted{, $\mathrm{Inf}$} dans $\mathcal{F}$ s'expriment par $\cup$, $\cap$ dans $\widetilde{\mathcal{F}}$.}

% The remaining lower block is cancelled by long diagonal strokes.
\deleted{\uline{At 3}}
\[
\deleted{F\mathrel{\unsure{\ll|\leq}}G\Longleftrightarrow\forall X\in\widetilde{F},\ \exists Y\in\widetilde{G}\text{ avec }X\ll Y}
\]
\deleted{\inserted{et si $\forall X\in\widetilde{F}$, $Y\in\widetilde{G}$,}}
\deleted{\uline{$X\cup Y$ existe} dans $\mathcal{F}$ (i.e. $\{X,Y\}$ majoré pour $\leq$), alors $F\cup G$ \unsure{.|:} existe dans $\mathcal{F}$.}

\deleted{On introduit relation dans $\deleted{\unsure{M}}$ $\mathcal{F}$}
\[
\deleted{\unsure{F}\mathrel{\gtrless}\unsure{G}\Longleftrightarrow\{\unsure{F},\unsure{G}\}\text{ majoré dans }\mathcal{F}}
\]
\deleted{\inserted{i.e. $F\cup G$ existe}}
\deleted{Par \uline{At 3}, il suffit de la connaître dans $M$.}
\scanpage{113}
\mspage{112}
% The top At 4 block is carried by the long curved arrow to the position before “+ les axiomes” below.
\deleted{\unsure{At}} \deleted{\uline{At 5} $F\mathrel{\unsure{|\circ|\,|\,|o|}}G$ implique que $F,G\in M$.}

\uline{At 3}\inserted{(spec)} \textcircled{a} Soient $F,G\in\unsure{\mathcal{F}|\mathscr{F}}$. Pour que $\{F,G\}$ majoré dans $\mathcal{F}^{\leq}$ i.e. $F\vee G$ existe, il f. et s. que
% slanted
\margin{Critère de compatibilité par strates}
\[
\begin{aligned}
&\deleted{\forall X\in\widetilde{F},\ Y\in\widetilde{G}}\\
&\deleted{\unsure{\forall}}X\in\widetilde{F}\setminus\widetilde{F}\cap\widetilde{G},\quad Y\in\widetilde{G}\setminus\widetilde{F}\cap\widetilde{G}\quad\deleted{\unsure{\text{on a}}}\Longrightarrow X\mathrel{|\circ|}Y
\end{aligned}
\]
\deleted{D'autre part, \unsure{la}}\par
\textcircled{b} \deleted{\unsure{de} $F,G$}
\[
F\mathrel{|\circ|}G\Longrightarrow F,G\in M,
\]
i.e. $|\circ|$ est en fait une relation dans $M$.

% The following insertion is physically at the top of the sheet; the arrow ends here.
\inserted{\uline{At 4}
\[
F\ll G\Longleftrightarrow\forall X\in\widetilde{F},\ \exists Y\in\widetilde{G},\ \text{avec }X\ll Y
\]
\deleted{\unsure{et}} Donc la connaissance de $\ll$ dans $\mathcal{F}$ équivaut à celle de $\ll$ dans $M$.}
% slanted
\margin{Critère de \unsure{raffinement} pour Ateliers}

$+$ les axiomes \uline{Mag} 1 -- \uline{Mag} 4 sur le $(M,\leq,\ll,|\circ|)$.

En termes de cette dernière structure, la donnée d'un \uline{atelier}\inserted{\uline{associé au magasin}} revient à la donnée de
\[
\widetilde{\mathcal{F}}\subset\unsure{P|\mathcal{P}}(M),
\]
satisfaisant les conditions

\uline{At M 1} $\forall\Phi\in\widetilde{\mathcal{F}}$, $\Phi$ est fermé pour $\leq$

\uline{At M 2} Si $\Phi\in\widetilde{\mathcal{F}}$, $\Psi\subset\Phi$, \deleted{alors}\inserted{est une} partie fermée de $\unsure{\varphi|\Phi}$, alors $\Psi\in\widetilde{\mathcal{F}}$

\uline{At M 3} \textcircled{a} $\displaystyle\bigcup_{\Phi\in\widetilde{\mathcal{F}}}\Phi=M$, et \textcircled{b} $\widetilde{\mathcal{F}}\ne\varnothing$ (ce dernier automatique si $M\ne\varnothing$), ou encore \textcircled{$b'$} $\varnothing\in\widetilde{\mathcal{F}}$.
\scanpage{114}
\mspage{113}
\uline{At M 4} Si $\Phi,\Psi\in\unsure{\widetilde{\mathcal{F}}|\widetilde{\mathscr{F}}}$, alors $\Phi\cup\Psi\in\widetilde{\mathcal{F}}$ ssi
\[
\begin{aligned}
&\forall X\in\Phi,\ Y\in\Psi,\ \deleted{\unsure{\text{on a}}\ X,Y\notin\Phi\cap\Psi,\ \text{on a }X\mathrel{\unsure{|\circ|\,|\,|o|}}Y.}\\
&\inserted{\text{on a }X\mathrel{\unsure{\gtrless|\lessgtr}}Y\ \text{i.e.}}\\
&\inserted{\forall X'\in\widetilde{X}\setminus\widetilde{X}\cap\widetilde{Y},\ Y'\in\widetilde{Y}\setminus\widetilde{X}\cap\widetilde{Y},\ \text{on a }\deleted{\unsure{X\mathrel{|\circ|}Y}}X'\mathrel{|\circ|}Y'.}
\end{aligned}
\]

L'atelier est dit \deleted{fini} ``à figures finies'' si les maquettes de ses figures sont finies. Cela revient \unsure{à|:} \unsure{se} donner $\deleted{\unsure{M}}$, le magasin correspondant, on trouve alors
\[
\widetilde{\mathcal{F}}=\{\deleted{\unsure{A}}\Phi\in\unsure{P|\mathcal{P}}_{f}(M)\mid X,Y\in\deleted{\unsure{A}}\Phi\Longrightarrow X\mathrel{\gtrless}Y\},
\]
mais \deleted{\unsure{est}} cela satisfait aux conditions d'un $M$-atelier ssi on a

\deleted{\unsure{At M}}\par
\uline{Mag fin} $\forall X\in M$, $\widetilde{X}$ est fini.

(Axiome des \uline{magasins} à \deleted{\unsure{multistrates finies}}\inserted{\uline{maquettes finies}}).

Au cours du travail, \inserted{dans cette petite \unsure{algèbre},} c'est finalement \uline{l'\unsure{étalier|atelier}} qui est au centre de l'attention constamment, alors que l'ensemble $\mathcal{F}$ de toutes les figures joue un rôle \unsure{strictement} accessoire.
